% Rewritten based on project.md
% LLNCS macro package for Springer Computer Science proceedings
% Version 2.21 of 2022/01/12
%
\documentclass[runningheads]{llncs}
%
\usepackage[T1]{fontenc}
\usepackage{graphicx}
\usepackage{amsmath}
\usepackage{booktabs}
\usepackage{hyperref}
\renewcommand\UrlFont{\color{blue}\rmfamily}
\urlstyle{rm}
%
\begin{document}
%
\title{Prediction of Water Saturation from Raw NMR Signals Using Deep Learning and Physics-Informed Approaches}

\titlerunning{Water Saturation Prediction from Raw NMR Signals}

\author{
Arick Jurdan dos Reis\inst{1} \and
Lorena Mamede Botelho\inst{1} \and
Claudio Miceli de Farias\inst{1} \and
Juan Pedro Perri Barreto\inst{1} \and
Marcio Mendes Taddei\inst{2}
}

\authorrunning{A. J. Reis et al.}

\institute{
Instituto T\'{e}rcio Pacitti de Aplica\c{c}\~{o}es e Pesquisas Computacionais (NCE),
Federal University of Rio de Janeiro (UFRJ), Rio de Janeiro, Brazil
\email{analytica@labnet.nce.ufrj.br}
\and
Brazilian Center for Research in Physics (CBPF),
Rio de Janeiro, Brazil
\email{mtaddei@cbpf.br}
\and
Energy Research Office (EPE),
Rio de Janeiro, Brazil
\email{arick.reis@epe.gov.br}
}

\maketitle
%
\begin{abstract}
Water saturation is a key petrophysical property for reservoir characterization and hydrocarbon volume estimation.
Conventional approaches rely on inverting Nuclear Magnetic Resonance (NMR) relaxation decays into $T_2$ distributions
and applying empirical cutoff values, a process that is both ill-posed and sensitive to modeling choices.
In this work we investigate whether machine-learning and deep-learning models can predict water saturation
directly from raw CPMG echo trains, bypassing the intermediate inversion step entirely.
We compare two input strategies---raw magnetization curves and inverted $T_2$ spectra---across a broad model
zoo that includes gradient-boosted trees (LightGBM, CatBoost), Random Forest, Logistic Regression, MLP,
CNNs, recurrent architectures (GRU, BiLSTM), Temporal Convolutional Networks, Temporal Fusion Transformers,
a Vanilla Transformer encoder, symbolic regression (PySR), and a physics-informed MLP.
A preprocessing pipeline combining ReLU activation, Discrete Wavelet Transform denoising,
normalization, and PCA-based dimensionality reduction is applied, with hyperparameters optimized via
the Tree-structured Parzen Estimator (TPE) in Optuna. Feature extraction from the time-series
signals is further explored using the TSFresh library at three complexity levels.
An ensemble of the best-performing conventional and physics-informed branches is proposed.
Models are evaluated on synthetic NMR data using MAE and RMSE.
Results indicate that the raw-curve approach, combined with appropriate preprocessing and
physics-informed regularization, can yield competitive saturation estimates while avoiding the
uncertainties inherent in $T_2$ inversion.

\keywords{NMR \and Water saturation \and Deep learning \and Physics-informed neural network \and Reservoir characterization}
\end{abstract}

%% ===================================================================
\section{Introduction}
\label{sec:intro}

Accurate determination of water saturation ($S_w$) is essential for estimating hydrocarbon volumes
and planning reservoir development. Nuclear Magnetic Resonance (NMR) logging is one of the primary
tools for this task: the Carr--Purcell--Meiboom--Gill (CPMG) sequence measures transverse
magnetization decays whose relaxation times encode information about pore-size distributions,
fluid types, and saturations~\cite{ref_coates1999}.

The standard workflow first inverts the raw echo train $M(t)$ into a $T_2$ relaxation-time
distribution using regularized inverse Laplace transforms (e.g., Tikhonov regularization).
A cutoff value is then applied to partition bound and free fluid, from which $S_w$ is inferred.
This two-step procedure, however, is inherently ill-posed: the inversion is sensitive to
regularization parameters and noise models, meaning that different implementations can yield
different $T_2$ distributions---and therefore different saturations---from the same raw data~\cite{ref_deepswirr}.
Setting the cutoff value introduces additional ambiguity, and several competing methods have been
proposed without a clear consensus~\cite{ref_liu2022,ref_solatpour2018}.

These difficulties are amplified in unconventional reservoirs such as shales, where strong
heterogeneity, multi-fluid signal overlap, and weak relaxation responses undermine conventional
cutoff- and template-based saturation methods~\cite{ref_li2026}.

Recent advances in deep learning (DL) offer an alternative path. DL models can learn complex
nonlinear mappings from high-dimensional inputs without explicit physical modeling and have already
been applied successfully to several petrophysical tasks: lithology classification~\cite{ref_valentin2019},
porosity and permeability prediction~\cite{ref_bom2021}, and subsurface
characterization~\cite{ref_kurup2006}. Notably, Reis et al.~\cite{ref_deepswirr} demonstrated
that a multi-scale convolutional network can estimate irreducible water saturation ($S_{w,irr}$)
directly from raw NMR echo trains with a correlation coefficient of approximately~0.9, bypassing
the $T_2$ inversion entirely. Similarly, Li et al.~\cite{ref_li2026} proposed a workflow that
combines unsupervised feature decoupling with a compact CNN to predict oil saturation from
$T_2$ spectra in shale reservoirs.

In this work, we extend the direct-prediction paradigm to water saturation by systematically
comparing a wide range of machine-learning and deep-learning architectures on synthetic NMR data.
We further propose a physics-informed MLP that incorporates known physical constraints into the
learning process, and combine it with conventional models in an ensemble pipeline.

%% ===================================================================
\section{Related Work}
\label{sec:related}

\subsection{DeepSWIrr}
Reis et al.~\cite{ref_deepswirr} estimated irreducible water saturation ($S_{w,irr}$)
directly from laboratory NMR echo trains using a multi-scale convolutional network.
Working with 821~carbonate-reservoir plugs from the Brazilian pre-salt, they preprocessed
the data by rotating channels to separate signal from noise and truncating each echo train
to 1000~points. Because labeled laboratory samples were scarce, a physically motivated
data-augmentation scheme was introduced: pairs or triplets of plugs are linearly combined
with a random weight~$w$, yielding synthetic magnetization curves
\begin{equation}
  m(w) = w\,m_1 + (1-w)\,m_2
  \label{eq:aug_m}
\end{equation}
and corresponding synthetic saturations
\begin{equation}
  S_w(w) =
  \frac{w\,S_{w,1}\,V_{B,1}\,\phi_1 + (1-w)\,S_{w,2}\,V_{B,2}\,\phi_2}
       {w\,V_{B,1}\,\phi_1 + (1-w)\,V_{B,2}\,\phi_2},
  \label{eq:aug_sw}
\end{equation}
where $V_B$ denotes bulk volume and $\phi$~porosity.
The model achieved a correlation of approximately~0.9 on a 50-plug blind test set,
with better performance on low-bulkiness samples.

\subsection{Direct Prediction of Oil Saturation from NMR $T_2$ Spectra}
Li et al.~\cite{ref_li2026} proposed a data-driven workflow for predicting oil saturation
from NMR~$T_2$ spectra in shale reservoirs. Their pipeline applies robust preprocessing,
dimensionality reduction, unsupervised clustering to decouple overlapping fluid signals,
and a lightweight CNN regressor. Validated on core--log paired data, the method showed
improved stability over conventional interpretation in complex intervals.

%% ===================================================================
\section{Methodology}
\label{sec:method}

\subsection{Data}
\label{sec:data}

We use synthetic NMR data generated following the procedure described in~\cite{ref_deepswirr},
under two key physical assumptions:
\begin{enumerate}
  \item Pore sizes follow a log-normal distribution.
  \item Each sample corresponds to a single pore system.
\end{enumerate}
% TODO: add details on synthetic data generation (number of samples, SNR levels, parameter ranges).

\subsection{Input Strategies}
\label{sec:input}

We investigate two complementary input representations:

\paragraph{Raw-curve approach.}
The input features~$\mathbf{X}$ are the temporal samples of the magnetization decay $M(t)$
(e.g., 2048~echoes from a CPMG sequence), and the target~$Y$ is the water
saturation~$S_w$ (or equivalently the oil saturation $S_o = 1 - S_w$).

\paragraph{Inverted-spectrum approach.}
The decay $M(t)$ is first inverted into a $T_2$ distribution using classical regularization
(e.g., Tikhonov). The input features become the amplitudes of the $T_2$ spectrum
(typically 64 or 128~bins). This representation is adopted in most of the existing literature.

\subsection{Preprocessing Pipeline}
\label{sec:preproc}

All input curves undergo the following sequential preprocessing steps (Fig.~\ref{fig:pipeline}):
\begin{enumerate}
  \item \textbf{ReLU clipping}: negative values are set to zero, removing unphysical amplitudes.
  \item \textbf{Discrete Wavelet Transform (DWT) denoising}: a wavelet-based smoothing is applied.
        The DWT is preferred over the traditional Savitzky--Golay filter because it better preserves
        sharp features in the decay while removing high-frequency noise.
  \item \textbf{Normalization}: curves are scaled to a standard range.
\end{enumerate}

\paragraph{Feature extraction with TSFresh.}
For tree-based and classical ML models, we additionally extract time-series features using the
TSFresh library~\cite{ref_tsfresh}. TSFresh applies statistical hypothesis tests to assess
feature relevance, reducing overfitting risk while scaling efficiently to multiple simultaneous
time series. We evaluate three extraction levels at two input lengths (700 and 2000~points):
\begin{itemize}
  \item \texttt{MinimalFCParameters}: fast and stable baseline features.
  \item \texttt{EfficientFCParameters}: a richer feature set with moderate computational cost.
  \item \texttt{ComprehensiveFCParameters}: an exhaustive but computationally expensive feature set.
\end{itemize}

\subsection{Models}
\label{sec:models}

We benchmark the following models, grouped by family:

\paragraph{Tree-based and classical ML.}
\begin{itemize}
  \item \textbf{LightGBM}~\cite{ref_lightgbm}: highly efficient gradient-boosted decision trees.
  \item \textbf{CatBoost}~\cite{ref_catboost}: gradient boosting with native categorical support.
  \item \textbf{Random Forest}: scikit-learn implementation.
  \item \textbf{Logistic Regression}: linear baseline (scikit-learn).
  \item \textbf{Symbolic Regression}: PySRRegressor via scikit-learn interface, seeking interpretable
        closed-form expressions.
\end{itemize}

\paragraph{Feedforward and physics-informed networks.}
\begin{itemize}
  \item \textbf{MLP}: multi-layer perceptron baseline (PyTorch).
  \item \textbf{Physics-Informed MLP}: an MLP augmented with physics-based loss terms
        encoding known NMR relaxation constraints (our approach).
\end{itemize}

\paragraph{Convolutional and temporal models.}
\begin{itemize}
  \item \textbf{CNN}: standard convolutional neural network (PyTorch).
  \item \textbf{Temporal Convolutional Network (TCN)}: dilated causal convolutions (PyTorch).
\end{itemize}

\paragraph{Recurrent models.}
\begin{itemize}
  \item \textbf{GRU}~\cite{ref_gru}: gated recurrent units.
  \item \textbf{BiLSTM}~\cite{ref_bilstm}: bidirectional long short-term memory, adapted
        for sequence-to-one regression.
\end{itemize}

\paragraph{Attention-based models.}
\begin{itemize}
  \item \textbf{Temporal Fusion Transformer (TFT)}~\cite{ref_tft}: originally designed for
        multi-horizon time-series forecasting; we adapt its architecture from sequence-to-sequence
        to sequence-to-scalar regression.
  \item \textbf{Vanilla Transformer}: encoder-only architecture.
\end{itemize}

\subsection{Hyperparameter Optimization}
\label{sec:optim}

All model hyperparameters are tuned using \textbf{Optuna}~\cite{ref_optuna}, with the default
Tree-structured Parzen Estimator (TPE) as the sampling strategy. The objective is to minimize
validation RMSE.

\subsection{Ensemble Strategy}
\label{sec:ensemble}

The final prediction is produced by an ensemble that combines the outputs of the best conventional
model branch (PCA $\rightarrow$ MLP) with the physics-informed MLP (PINN) branch, as illustrated
in Fig.~\ref{fig:pipeline}.

\begin{figure}[t]
  \centering
  % TODO: replace with actual pipeline figure
  % The pipeline follows: Data -> ReLU -> DWT -> Norm -> {PCA -> MLP, PINN} -> Ensemble -> Output
  \fbox{\parbox{0.9\textwidth}{\centering\vspace{2em}%
    \textbf{Pipeline diagram placeholder}\\[0.5em]
    Data $\rightarrow$ ReLU $\rightarrow$ DWT $\rightarrow$ Norm
    $\rightarrow$ \{PCA $\rightarrow$ MLP,\; PINN\}
    $\rightarrow$ Ensemble $\rightarrow$ Output
  \vspace{2em}}}
  \caption{Overview of the proposed preprocessing and modeling pipeline. Raw NMR decay curves
  are preprocessed (ReLU clipping, DWT denoising, normalization) and then fed into two parallel
  branches: a PCA-reduced MLP and a physics-informed neural network (PINN). The outputs are
  combined in an ensemble to produce the final water-saturation estimate.}
  \label{fig:pipeline}
\end{figure}

\subsection{Evaluation Metrics}
\label{sec:metrics}

All models are evaluated using:
\begin{itemize}
  \item \textbf{Mean Absolute Error (MAE)}: $\text{MAE} = \frac{1}{n}\sum_{i=1}^{n}|y_i - \hat{y}_i|$
  \item \textbf{Root Mean Squared Error (RMSE)}: $\text{RMSE} = \sqrt{\frac{1}{n}\sum_{i=1}^{n}(y_i - \hat{y}_i)^2}$
\end{itemize}

%% ===================================================================
\section{Results}
\label{sec:results}

% TODO: Insert experimental results here.
% No experimental results are available yet. Tables and figures comparing model
% performance (MAE, RMSE) across all architectures and input strategies should be
% added once experiments are complete.

\textit{Experimental results are pending. This section will report MAE and RMSE for all models
under both input strategies (raw curve and inverted spectrum), together with ablation studies
on preprocessing steps and feature-extraction configurations.}

%% ===================================================================
\section{Conclusion}
\label{sec:conclusion}

% TODO: Write conclusions after results are available.

\textit{Conclusions will be drawn once experimental evaluation is complete. We expect to
demonstrate that the raw-curve approach, augmented with physics-informed regularization and
ensemble strategies, offers a viable alternative to the conventional $T_2$-inversion pipeline
for water-saturation prediction from NMR data.}

%% ===================================================================
\begin{credits}
\subsubsection{\ackname}
The authors thank the Brazilian Center for Research in Physics (CBPF),
Petrobras, and the LabNet laboratory at UFRJ for supporting this research.

\subsubsection{\discintname}
The authors have no competing interests to declare that are relevant to the
content of this article.
\end{credits}

%% ===================================================================
% Bibliography
%
\bibliographystyle{splncs04}
%
\begin{thebibliography}{15}

\bibitem{ref_coates1999}
Coates, G.R., Xiao, L., Prammer, M.G.:
NMR Logging: Principles and Applications.
Halliburton Energy Services, Houston (1999)

\bibitem{ref_deepswirr}
Reis, A.J., et al.:
DeepSWIrr: Determining irreducible water saturation from raw Nuclear Magnetic Resonance
decays using deep learning.
% TODO: add journal/conference, volume, pages, year

\bibitem{ref_liu2022}
Liu, X., et al.:
T2 cutoff value determination methods for NMR logging.
% TODO: complete citation (2022)

\bibitem{ref_solatpour2018}
Solatpour, A., et al.:
NMR T2 cutoff determination methods.
% TODO: complete citation (2018)

\bibitem{ref_li2026}
Li, Q., Wang, C., Ge, X., Chi, R., Wang, Y., Zhang, W., Miao, Q., Zhang, J.:
Direct Prediction of Oil Saturation from NMR $T_2$ Spectra Using Unsupervised Feature
Decoupling and Deep Learning.
In: EGU General Assembly 2026, Vienna, Austria, EGU26-3276 (2026).
\doi{10.5194/egusphere-egu26-3276}

\bibitem{ref_valentin2019}
Valent\'{i}n, M.B., et al.:
A deep learning approach for lithology classification using well log data.
% TODO: complete citation (2019)

\bibitem{ref_bom2021}
Bom, C.R., et al.:
Deep learning for porosity and permeability prediction from well logs.
% TODO: complete citation (2021)

\bibitem{ref_kurup2006}
Kurup, P.U., Griffin, E.P.:
Prediction of soil composition from CPT data using general regression neural network.
Journal of Computing in Civil Engineering \textbf{20}(4), 281--289 (2006)

\bibitem{ref_tsfresh}
Christ, M., Braun, N., Neuffer, J., Kempa-Liehr, A.W.:
Time Series FeatuRe Extraction on basis of Scalable Hypothesis tests (tsfresh -- A Python package).
Neurocomputing \textbf{307}, 72--77 (2018)

\bibitem{ref_lightgbm}
Ke, G., Meng, Q., Finley, T., Wang, T., Chen, W., Ma, W., Ye, Q., Liu, T.-Y.:
LightGBM: A Highly Efficient Gradient Boosting Decision Tree.
In: Advances in Neural Information Processing Systems~30, pp.~3146--3154 (2017)

\bibitem{ref_catboost}
Prokhorenkova, L., Gusev, G., Vorobev, A., Dorogush, A.V., Gulin, A.:
CatBoost: unbiased boosting with categorical features.
In: Advances in Neural Information Processing Systems~31, pp.~6638--6648 (2018)

\bibitem{ref_gru}
Chung, J., Gulcehre, C., Cho, K., Bengio, Y.:
Empirical Evaluation of Gated Recurrent Neural Networks on Sequence Modeling.
arXiv preprint arXiv:1412.3555 (2014)

\bibitem{ref_bilstm}
MathWorks:
Sequence-to-One Regression Using Deep Learning.
MATLAB \& Simulink Documentation.
\url{https://www.mathworks.com/help/deeplearning/ug/sequence-to-one-regression-using-deep-learning.html}

\bibitem{ref_tft}
Lim, B., Ar\'{i}k, S.\"{O}., Loeff, N., Pfister, T.:
Temporal Fusion Transformers for Interpretable Multi-horizon Time Series Forecasting.
International Journal of Forecasting \textbf{37}(4), 1748--1764 (2021)

\bibitem{ref_optuna}
Akiba, T., Sano, S., Yanase, T., Ohta, T., Koyama, M.:
Optuna: A Next-generation Hyperparameter Optimization Framework.
In: Proceedings of the 25th ACM SIGKDD International Conference on Knowledge Discovery
\& Data Mining, pp.~2623--2631 (2019)

\end{thebibliography}
\end{document}
